% !TEX program = xelatex
% ============================================================
%  خطة تحسين SEO + إصلاح اسم الموقع والأيقونة في نتائج جوجل
%  الاستخدام على Overleaf:
%    1) أنشئ مشروعًا جديدًا فارغًا (Blank Project)
%    2) الصق محتوى هذا الملف في main.tex
%    3) Menu ← Compiler ← اختر: XeLaTeX   (مهم للعربية)
%    4) اضغط Recompile ← ستحصل على PDF
% ============================================================
\documentclass[11pt]{article}

\usepackage{polyglossia}
\setmainlanguage{arabic}
\setotherlanguage{english}
\newfontfamily\arabicfont[Script=Arabic,Scale=1.0]{Amiri}
\newfontfamily\arabicfontsf[Script=Arabic,Scale=1.0]{Amiri}
\setmonofont{DejaVu Sans Mono}[Scale=0.8]

\usepackage{geometry}
\geometry{a4paper, margin=2.2cm}

\usepackage{xcolor}
\usepackage{graphicx}
\usepackage{booktabs}
\usepackage{longtable}
\usepackage{array}
\usepackage{enumitem}
\usepackage{titlesec}
\usepackage{fancyhdr}
\usepackage{hyperref}
\usepackage{tcolorbox}
\tcbuselibrary{skins,breakable}
\usepackage{fontawesome5}

% ── الألوان (هوية يَسِّر) ──
\definecolor{gold}{HTML}{B8933B}
\definecolor{darkbg}{HTML}{101414}
\definecolor{okgreen}{HTML}{27AE60}
\definecolor{warnred}{HTML}{C0392B}
\definecolor{softgray}{HTML}{F4F4F4}
\definecolor{midgray}{HTML}{7A7A7A}
\definecolor{codebg}{HTML}{1A2020}

% ── الترويسة والتذييل ──
\pagestyle{fancy}
\fancyhf{}
\renewcommand{\headrulewidth}{0.4pt}
\renewcommand{\footrulewidth}{0.4pt}
\fancyhead[R]{\small\color{gold}خطة تحسين SEO — منصة يَسِّر}
\fancyhead[L]{\small\color{midgray}Search Optimization Roadmap}
\fancyfoot[C]{\small\thepage}

% ── عناوين ──
\titleformat{\section}{\Large\bfseries\color{gold}}{\thesection}{0.6em}{}
\titleformat{\subsection}{\large\bfseries\color{darkbg}}{\thesubsection}{0.5em}{}
\titleformat{\subsubsection}{\normalsize\bfseries\color{midgray}}{\thesubsubsection}{0.4em}{}

\setlist[itemize]{rightmargin=0pt, itemsep=2pt, topsep=3pt}
\setlist[enumerate]{rightmargin=0pt, itemsep=2pt, topsep=3pt}

\hypersetup{
  colorlinks=true, linkcolor=gold, urlcolor=gold,
  pdftitle={خطة تحسين SEO - منصة يَسِّر}, pdfauthor={فريق يَسِّر}
}

% ── صناديق ──
\newtcolorbox{infobox}[1][]{colback=softgray,colframe=gold,boxrule=0.6pt,arc=3pt,left=8pt,right=8pt,top=6pt,bottom=6pt,breakable,#1}
\newtcolorbox{riskbox}[1][]{colback=red!4,colframe=warnred,boxrule=0.6pt,arc=3pt,left=8pt,right=8pt,top=6pt,bottom=6pt,breakable,#1}
\newtcolorbox{okbox}[1][]{colback=green!4,colframe=okgreen,boxrule=0.6pt,arc=3pt,left=8pt,right=8pt,top=6pt,bottom=6pt,breakable,#1}

\newcommand{\code}[1]{\textcolor{darkbg}{\texttt{\small #1}}}

\begin{document}

% ===================== الغلاف =====================
\begin{titlepage}
\centering
\vspace*{2.5cm}
{\Huge\bfseries\color{gold} يَسِّر}\\[0.4cm]
{\large\color{midgray} منصة تعليم الفقه الإسلامي}\\[2.5cm]

{\Huge\bfseries خطة تحسين محركات البحث}\\[0.8cm]
{\LARGE\bfseries SEO — اسم الموقع، الأيقونة، والظهور}\\[1.2cm]

\rule{0.7\textwidth}{0.5pt}\\[0.8cm]
{\large Search Engine Optimization Roadmap}\\[0.8cm]
\rule{0.7\textwidth}{0.5pt}\\[2.5cm]

\begin{tabular}{r l}
\textbf{الهدف الأول:}   & اسم «يَسِّر» في جوجل بدل «الباحث الفقهي» \\
\textbf{الهدف الثاني:}  & أيقونة الموقع في نتائج البحث \\
\textbf{الهدف الثالث:}  & تحسين عام للفهرسة والثِقة \\
\textbf{المدة التقنية:} & ≈ ٣ أيام عمل \\
\textbf{زمن الظهور:}    & ٢–٦ أسابيع (فهرسة جوجل) \\
\textbf{النسخة:}        & 1.0 \\
\end{tabular}

\vfill
{\small\color{midgray} وثيقة تخطيط تقني — \today}
\end{titlepage}

\newpage
\tableofcontents
\newpage

% ===================== الملخص =====================
\section{الملخص التنفيذي}

\begin{infobox}
يهدف هذا المستند إلى حل ثلاث مشكلات ظهور:
\begin{enumerate}
\item ظهور اسم قديم «الباحث الفقهي» في نتائج جوجل بدل الاسم الرسمي «يَسِّر».
\item عدم ظهور أيقونة الموقع بجانب النتيجة.
\item ضعف الإشارات البنيوية (\code{JSON-LD}) التي تعتمد عليها جوجل لتسمية الموقع.
\end{enumerate}
الحل يقوم على إشارات قاطعة (\code{WebSite} + \code{Organization} JSON-LD)، وأيقونات قياسية (\code{favicon.ico} + PNG مربّعة)، ثم إعادة فهرسة موجّهة عبر Google Search Console.
\end{infobox}

\subsection*{التشخيص — سبب كل مشكلة}
\begin{center}
\renewcommand{\arraystretch}{1.35}
\begin{tabular}{@{}p{4.5cm} p{9.5cm}@{}}
\toprule
\textbf{المطلوب} & \textbf{السبب الجذري الحالي} \\
\midrule
اسم «يَسِّر» في جوجل & لا يوجد \code{WebSite} JSON-LD؛ الاعتماد على \code{og:site\_name} وحده ضعيف. الاسم القديم بقايا فهرسة لم تُحدَّث \\
أيقونة الموقع في النتائج & فقط \code{favicon.svg}؛ جوجل يطلب \code{favicon.ico} في الجذر + أيقونة مربّعة ≥ 48px \\
تحسين SEO & \code{<title>} قصير جدًا؛ لا \code{SearchAction}؛ لا \code{BreadcrumbList} \\
\bottomrule
\end{tabular}
\end{center}

% ===================== المراحل =====================
\section{المراحل}

\subsection{المرحلة ١ — إشارة اسم الموقع (الأهم) — نصف يوم}
الهدف: إخبار جوجل قطعيًا أن اسم الموقع «يَسِّر».

\begin{tabular}{@{}c p{6.5cm} p{5cm}@{}}
\toprule
\textbf{\#} & \textbf{المهمة} & \textbf{التفصيل} \\
\midrule
١.١ & \code{WebSite} JSON-LD في \code{index.html} & \code{name}, \code{alternateName}, \code{url}, \code{inLanguage} \\
١.٢ & \code{Organization} JSON-LD & \code{name}, \code{logo}, \code{sameAs}, \code{url} \\
١.٣ & \code{SearchAction} (روابط البحث) & \code{potentialAction} ← مسار البحث \code{/\#/lessons} \\
١.٤ & تحديث \code{og:site\_name} و\code{og:title} & «يَسِّر — منصة تعليم الفقه الإسلامي» \\
١.٥ & تصحيح \code{<title>} الرئيسي & «يَسِّر — منصة تعليم الفقه الإسلامي الميسّر» \\
\bottomrule
\end{tabular}

\begin{infobox}
\textbf{نموذج الـJSON-LD المطلوب:}
\begin{center}
\begin{tcolorbox}[colback=codebg, colframe=codebg, coltext=white, arc=3pt, left=10pt, right=10pt, width=\linewidth]
\ttfamily\footnotesize\raggedleft
\{\\
\hspace*{1em}"@context": "https://schema.org",\\
\hspace*{1em}"@graph": [\\
\hspace*{2em}\{ "@type": "WebSite", "name": "يَسِّر",\\
\hspace*{3em}"alternateName": ["يسر", "منصة يَسِّر للفقه"],\\
\hspace*{3em}"url": "https://fqh.me/", "inLanguage": "ar",\\
\hspace*{3em}"potentialAction": \{\\
\hspace*{4em}"@type": "SearchAction",\\
\hspace*{4em}"target": "https://fqh.me/\#/lessons?q=\{search\_term\_string\}" \} \},\\
\hspace*{2em}\{ "@type": "Organization", "name": "يَسِّر",\\
\hspace*{3em}"url": "https://fqh.me/",\\
\hspace*{3em}"logo": "https://fqh.me/icons/icon-512.png" \}\\
\hspace*{1em}]\\
\}
\end{tcolorbox}
\end{center}
\end{infobox}

\subsection{المرحلة ٢ — الأيقونة في نتائج جوجل — نصف يوم}
الهدف: ظهور أيقونة «يَسِّر» بجانب نتيجة البحث.

\begin{tabular}{@{}c p{6.5cm} p{5cm}@{}}
\toprule
\textbf{\#} & \textbf{المهمة} & \textbf{التفصيل} \\
\midrule
٢.١ & توليد أيقونات PNG من الشعار & 48، 96، 192، 512 \\
٢.٢ & إنشاء \code{favicon.ico} متعدد الأحجام & 16 + 32 + 48 داخل ملف واحد \\
٢.٣ & روابط الأيقونات في \code{index.html} & \code{rel=icon sizes=48x48} + apple-touch \\
٢.٤ & أيقونة البار في جوجل & أيقونة مربّعة ≥ 48px + \code{favicon.ico} في الجذر \\
٢.٥ & إصلاح \code{favicon.svg} & حاليًا يضمّن PNG base64 (12KB) — يُستحسن SVG حقيقي \\
\bottomrule
\end{tabular}

\begin{riskbox}
أيقونة نتيجة البحث تأخذ أسابيع لتحديثها، ويجب أن يكون \code{favicon.ico} في \textbf{جذر الموقع} لا في مسار فرعي.
\end{riskbox}

\subsection{المرحلة ٣ — إعادة الفهرسة وتحديث الاسم — يوم + انتظار}
الهدف: محو «الباحث الفقهي» القديم من نتائج جوجل.

\begin{tabular}{@{}c p{6.5cm} p{5cm}@{}}
\toprule
\textbf{\#} & \textbf{المهمة} & \textbf{التفصيل} \\
\midrule
٣.١ & Google Search Console & تأكيد الملكية ← طلب إعادة فهرسة للرئيسية \\
٣.٢ & URL Inspection في GSC & فحص تسمية الموقع بعد النشر \\
٣.٣ & تسمية الموقع (Site name) & تُشتق من JSON-LD أو \code{og:site\_name} \\
٣.٤ & تأكيد \code{robots.txt} + sitemap & إعادة إرسال \code{sitemap.xml} في GSC \\
٣.٥ & تحديث \code{lastmod} في sitemap & التأكد من التحديث بعد النشر \\
٣.٦ & Bing Webmaster & اختياري: إبلاغ Bing أيضًا \\
\bottomrule
\end{tabular}

\subsection{المرحلة ٤ — تحسينات SEO أعمق — يوم}
\begin{tabular}{@{}c p{9.5cm}@{}}
\toprule
\textbf{\#} & \textbf{المهمة (والأثر)} \\
\midrule
٤.١ & \code{BreadcrumbList} JSON-LD في صفحات المسائل ← مسار التنقل في النتائج \\
٤.٢ & إثراء النتائج بـFAQ/Topic schema (اختياري) \\
٤.٣ & تحسين الوصف الميتا (١٢٠–١٦٠ حرفًا) \\
٤.٤ & \code{<meta name=robots content=index,follow>} صريح \\
٤.٥ & تحسين Core Web Vitals (LCP/CLS/INP) — مؤثر في الترتيب \\
٤.٦ & فقرة نصّية تعريفية في الرئيسية \\
٤.٧ & تحديث \code{check:seo.mjs} لفحص JSON-LD والاسم \\
\bottomrule
\end{tabular}

\subsection{المرحلة ٥ — التحقق والقياس — نصف يوم}
\begin{tabular}{@{}c p{6cm} p{6cm}@{}}
\toprule
\textbf{\#} & \textbf{الفحص} & \textbf{الأداة} \\
\midrule
٥.١ & فحص JSON-LD & validator.schema.org \\
٥.٢ & فحص النتائج المثراة & Rich Results Test \\
٥.٣ & فحص اسم الموقع & GSC ← تسمية الموقع \\
٥.٤ & فحص الأيقونة & RealFaviconGenerator \\
٥.٥ & Lighthouse SEO & ≥ ٩٥ \\
٥.٦ & فحص البناء & \code{npm run check:seo} \\
\bottomrule
\end{tabular}

% ===================== الحرجة =====================
\section{ملاحظات حرجة}
\begin{riskbox}
\begin{itemize}
\item \code{notranslate} موجود في \code{index.html} — جيد؛ يمنع الترجمة الآلية التي تشوّه الاسم.
\item اسم «الباحث الفقهي» فُهرس قبل تغيير الاسم؛ جوجل يحتاج \textbf{٢–٦ أسابيع} للتحديث بعد إعادة الفهرسة. \textbf{لا يمكن تسريعه قسرًا}.
\item الاسم في \code{<title>} هو الأقوى — حاليًا «يَسِّر» فقط، فيجب أن يكون صريحًا.
\item تأكد أن \code{fqh.me} نهائي (لا \code{www} يُحوّل) — الحالي صحيح.
\end{itemize}
\end{riskbox}

% ===================== الجدول =====================
\section{الجدول الزمني}
\begin{center}
\renewcommand{\arraystretch}{1.3}
\begin{tabular}{@{}p{7cm} c c@{}}
\toprule
\textbf{المرحلة} & \textbf{المدة} & \textbf{خارجي} \\
\midrule
١ — JSON-LD الاسم & ٠.٥ يوم & — \\
٢ — الأيقونات & ٠.٥ يوم & — \\
٣ — إعادة الفهرسة & ٠.٥ يوم & \textcolor{warnred}{٢–٦ أسابيع} \\
٤ — SEO أعمق & ١ يوم & — \\
٥ — التحقق & ٠.٥ يوم & — \\
\midrule
\textbf{الإجمالي} & \textbf{≈ ٣ أيام} & — \\
\bottomrule
\end{tabular}
\end{center}

% ===================== Done =====================
\section{معايير الإنجاز (Definition of Done)}
\begin{okbox}
\begin{itemize}
\item[\faCheck] \code{index.html} يضم JSON-LD: \code{WebSite} + \code{Organization} + \code{SearchAction}.
\item[\faCheck] أيقونات: \code{favicon.ico} + أيقونات PNG بمقاسات ٤٨ و٩٦ و١٩٢ و٥١٢.
\item[\faCheck] \code{og:site\_name} و\code{<title>} محدّثتان رسميًا.
\item[\faCheck] \code{generate-seo.mjs} يولّد \code{BreadcrumbList} والأيقونات.
\item[\faCheck] GSC: إعادة فهرسة + تأكيد اسم الموقع.
\item[\faCheck] \code{check:seo.mjs} يفحص الاسم وJSON-LD (يفشل عند الكسر).
\end{itemize}
\end{okbox}

% ===================== ملحق =====================
\section{ملحق — هيكل الملفات بعد التنفيذ}
\begin{tcolorbox}[colback=darkbg, colframe=darkbg, coltext=white, arc=3pt, left=8pt, right=8pt]
\ttfamily\small\raggedleft
index.html\hspace{1em}\textcolor{okgreen}{مُعدّل — JSON-LD + أيقونات}\\
public/\\
\hspace*{1em}favicon.ico\hspace{1em}\textcolor{gold}{جديد}\\
\hspace*{1em}favicon.svg\hspace{1em}\textcolor{okgreen}{مُحسَّن}\\
\hspace*{1em}icons/\\
\hspace*{2em}icon-48.png\hspace{1em}\textcolor{gold}{جديد}\\
\hspace*{2em}icon-96.png\hspace{1em}\textcolor{gold}{جديد}\\
\hspace*{2em}icon-192.png\hspace{1em}\textcolor{gold}{جديد}\\
\hspace*{2em}icon-512.png\hspace{1em}\textcolor{gold}{جديد}\\
scripts/\\
\hspace*{1em}generate-seo.mjs\hspace{1em}\textcolor{okgreen}{مُعدّل — BreadcrumbList}\\
\hspace*{1em}check-seo.mjs\hspace{1em}\textcolor{okgreen}{مُعدّل — فحص الاسم}
\end{tcolorbox}

\vfill
\begin{center}
\rule{0.5\textwidth}{0.4pt}\\[0.4cm]
{\small\color{midgray} نهاية الوثيقة — منصة يَسِّر}
\end{center}

\end{document}
